\section{Introduction}
\label{sec:introduction}

ADMET models help prioritize compounds for synthesis and experimental testing. A useful model must
therefore do more than make accurate predictions: it should also indicate when those predictions
are uncertain. Reliable uncertainty estimates can help decide which compounds to test, where more
data are needed, and when a prediction should not drive an experimental decision
\citep{hirschfeld2020uncertainty}.

ADMET datasets are often small, and a more complicated model is not necessarily better. Fingerprint
and descriptor models remain competitive across many small molecular datasets
\citep{notwell2023maplight,praski2025benchmarking}, while predictability differs substantially
across assays \citep{aleksic2022predictability}. Uncertainty quantification
(UQ) methods may likewise perform differently from one endpoint to another.

The comparison itself also varies. Chemical partitions, model initialization, and the data used for
tuning can change the observed differences between both point predictors and UQ methods
\citep{bouthillier2021accounting,ash2025practically}. This matters for the TDC ADMET benchmark, which
contains 22 public prediction tasks with endpoint-specific performance measures and a fixed scaffold
test split \citep{huang2021tdc}. After our structure curation, the datasets contain 475 to 13,094
compounds. An audit of leading
submissions found reproducibility failures, data leakage, and model choices influenced by the public
test sets \citep{koleiev2026enduser}. We therefore ask which methods perform consistently across
datasets, which work best on particular endpoints, and whether those conclusions survive changes in
the evaluation design.

We study eleven calibrated pipelines based on LightGBM and single or ensemble neural models on all
22 datasets, crossing ten chemical partitions, five model seeds, and five outer folds. This grid
supports comparisons across datasets and on each dataset's original score scale, followed by
smaller-design stability analyses.

Our contributions are:
\begin{itemize}
  \item We compare eleven calibrated UQ pipelines on all 22 TDC ADMET datasets using nested
  cross-validation across multiple chemical partitions and model seeds.
  \item We show that ensemble MC dropout is the most consistent method across datasets, while
  LightGBM wins many individual endpoints and performs poorly on others.
  \item We examine how reduced split and tuning designs affect both the apparent leader and the
  differences supported by the statistical analysis.
\end{itemize}
